% =============================================================
% review.tex — Звіт з практики за темою БКР
% Тема: Візуальне моделювання роботи обчислювача ШПФ
%       під керуванням потоку даних
% НУ «Львівська політехніка», каф. АСУ, 2026
% =============================================================

\documentclass[14pt,a4paper]{extarticle}

% --- Мова і кодування ---
\usepackage[utf8]{inputenc}
\usepackage[ukrainian,english]{babel}

% --- Шрифт Times New Roman ---
\usepackage{tempora}
\usepackage[T1,T2A]{fontenc}

% --- Поля (вимоги кафедри АСУ) ---
\usepackage[
  left=25mm,
  right=10mm,
  top=20mm,
  bottom=20mm
]{geometry}

% --- Міжрядковий інтервал 1.5 ---
\usepackage{setspace}
\onehalfspacing

% --- Математика ---
\usepackage{amsmath}
\usepackage{amssymb}
\usepackage{amsfonts}
\usepackage{icomma}

% --- Таблиці та зображення ---
\usepackage{graphicx}
\usepackage{float}
\usepackage{booktabs}
\usepackage{array}
\usepackage{longtable}

% --- Відступи і параграфи ---
\usepackage{indentfirst}
\usepackage{microtype}
\emergencystretch=3em
\setlength{\parindent}{1.25cm}

% --- Переноси ---
\usepackage[htt]{hyphenat}

% --- Налаштування заголовків ---
\usepackage{titlesec}
\titleformat{\section}[block]
  {\normalfont\large\bfseries\centering}
  {\thesection.}{0.5em}{}
\titleformat{\subsection}[block]
  {\normalfont\normalsize\bfseries}
  {\thesubsection}{0.5em}{}
\titleformat{\subsubsection}[block]
  {\normalfont\normalsize\bfseries}
  {\thesubsubsection}{0.5em}{}

% --- Гіперпосилання ---
\usepackage{hyperref}
\hypersetup{
  colorlinks=true,
  linkcolor=black,
  citecolor=black,
  urlcolor=blue,
  pdfauthor={Тесля М.},
  pdftitle={Звіт з практики за темою БКР}
}

\begin{document}

% =============================================================
\begin{center}
  \large\bfseries
  ЗВІТ З ПРАКТИКИ ЗА ТЕМОЮ \\[2pt]
  БАКАЛАВРСЬКОЇ КВАЛІФІКАЦІЙНОЇ РОБОТИ \\[6pt]
  \normalsize\normalfont
  «Візуальне моделювання роботи обчислювача ШПФ \\
  під керуванням потоку даних»
\end{center}

\vspace{1em}

% =============================================================
\section{Робота над темою БКР}

% ─────────────────────────────────────────────────
\subsection{Постановка задачі та визначення основних функціональних
вимог і завдань роботи}

Темою бакалаврської кваліфікаційної роботи (БКР) є розробка
веб-орієнтованого програмного симулятора потокового обчислювача
16-точкового швидкого перетворення Фур'є (ШПФ) з основою~4
($\text{radix-}4$, $4 \times 4$). Актуальність роботи зумовлена
фундаментальною роллю ШПФ у задачах цифрової обробки сигналів
(OFDM-модуляція у LTE, 5G, Wi-Fi, обробка мови, радіолокація,
спектральний аналіз) і одночасно --- наявністю системного розриву
між математичним описом алгоритму, його формальною моделлю потоку
даних та виробничими реалізаціями~[1][2].

\textbf{Формальна постановка задачі.}
Нехай задано комплекснозначну вхідну послідовність
$\{x_n\}_{n=0}^{N-1}$ довжиною $N = 16$, подану у подвійній
точності формату IEEE~754. Потрібно розробити програмний симулятор,
який обчислює дискретне перетворення Фур'є (ДПФ)

\begin{equation*}
  X_k = \sum_{n=0}^{N-1} x_n \, e^{-j 2\pi n k / N},
  \quad k = 0, 1, \ldots, N-1,
\end{equation*}

за алгоритмом radix-4 з двовимірною $4 \times 4$ декомпозицією,
формалізує процес обчислення у вигляді орієнтованого графа потоку
даних (DFG, англ.\ \textit{Data Flow Graph}) і забезпечує його
інтерактивну візуалізацію з анімованим рухом токенів, метриками та
числовою верифікацією~[3][4].

\textbf{Вхідні дані симулятора:} 16~комплексних відліків сигналу
(задаються користувачем або генеруються вбудованим джерелом ---
синусоїда, одиничний імпульс, білий шум); параметри симуляції
(затримки вузлів та дуг, множник швидкості); команди керування
(пуск, пауза, покроковий режим, скидання).

\textbf{Вихідні дані:} спектр $\{X_k\}_{k=0}^{15}$, що відображається
як амплітуда $|X_k|$ і фаза $\arg X_k$; журнал подій
(\texttt{onToken}, \texttt{onFire}, \texttt{onOutput}); метрики
(пропускна здатність, наскрізна затримка, розмір черг, теплова карта
активності вузлів); звіт про верифікацію --- максимальне відхилення
від еталонного ДПФ.

\textbf{Функціональні вимоги.}
На основі аналізу предметної області визначено перелік із дев'яти
функціональних вимог F1--F9 (табл.~\ref{tab:func_req}).

\begin{table}[H]
  \caption{Функціональні вимоги до симулятора}
  \label{tab:func_req}
  \centering
  \small
  \begin{tabular}{p{1.0cm}p{4.0cm}p{8.5cm}}
    \toprule
    ID & Назва & Опис \\
    \midrule
    F1 & Побудова DFG               & Автоматична генерація графа потоку даних для ШПФ radix-4 $4{\times}4$ при $N{=}16$ \\
    F2 & Симуляція виконання        & Крокування часу з доставкою токенів, перевіркою правила активації та запуском вузлів \\
    F3 & Введення сигналу           & Вибір типу вхідного сигналу та його параметрів \\
    F4 & Керування симуляцією       & Пуск, пауза, покроковий режим, скидання, зміна швидкості \\
    F5 & Візуалізація токенів       & Анімоване переміщення токенів по дугах із відображенням значень \\
    F6 & Відображення метрик        & Throughput, latency, queue size, теплова карта активності \\
    F7 & Спектральний вихід         & Гістограма $|X_k|$ і $\arg X_k$ \\
    F8 & Верифікація результату     & Порівняння з еталонним ДПФ і виведення максимального відхилення \\
    F9 & Налаштування параметрів    & Редагування \texttt{latency}, \texttt{delay} без перезавантаження сторінки \\
    \bottomrule
  \end{tabular}
\end{table}

\textbf{Нефункціональні вимоги.}
Якісні характеристики системи регламентуються вісьмома
нефункціональними вимогами NF1--NF8 у відповідності до стандарту
ISO/IEC~25010:2011: точність обчислень
$\max_k |X_k^{\mathrm{sim}} - X_k^{\mathrm{ref}}| \le 10^{-9}$ (NF1);
час відгуку UI $\le 100$~мс (NF2); стабільні 60~fps при швидкості
симуляції $\le 5\times$ (NF3); портативність у Chrome~$\ge$~90,
Firefox~$\ge$~88, Edge~$\ge$~90 (NF4); детермінованість виконання
(NF5); відкрита ліцензія MIT/Apache~2.0 (NF6); робота у браузері без
серверної частини (NF7); зручність --- основні операції доступні не
глибше двох кліків від головного екрана (NF8).

\textbf{Завдання роботи.}
Виходячи з постановки задачі та переліку вимог, сформульовано такі
основні завдання:
\begin{enumerate}
  \item проаналізувати предметну область, провести огляд літератури
        та існуючих рішень для обчислення ШПФ;
  \item виконати системний аналіз задачі (зацікавлені сторони, цілі,
        обмеження, критерії успіху);
  \item формалізувати DFG-граф 16-точкового ШПФ radix-4 з вузлами
        типів \texttt{source}, \texttt{dft4}, \texttt{twiddle},
        \texttt{sink} та дугами з параметричними затримками;
  \item обґрунтувати вибір технологій реалізації;
  \item розробити структуру системи (архітектура C4, структури даних,
        механізми обміну даними, інтерфейс користувача);
  \item реалізувати ядро симуляції з правилом активації токенів,
        FIFO-чергами та керованим логічним часом;
  \item забезпечити інтерактивну візуалізацію руху токенів і
        спектральний вихід;
  \item виконати числову верифікацію результатів шляхом порівняння
        з еталонним ДПФ.
\end{enumerate}

% ─────────────────────────────────────────────────
\subsection{Огляд літератури та аналогів}

Аналіз літературних джерел і наявних програмних рішень за останні
10~років показав, що простір сучасних засобів обчислення ШПФ
поділяється на три сегменти: програмні бібліотеки, апаратні та
GPU-прискорювачі, системи автоматичної генерації коду. Окремо
розглянуто методи візуального моделювання обчислювальних систем.

\subsubsection*{Програмні бібліотеки обчислення ШПФ}

\textbf{FFTW} (Fastest Fourier Transform in the West, MIT, 1997) ---
бібліотека з адаптивним планувальником, що методом динамічного
програмування добирає оптимальну рекурсивну факторизацію під
конкретну апаратну платформу. Архітектура побудована на
кеш-незалежній (\textit{cache-oblivious}) рекурсії «розділяй і
володарюй»~[8].

\textbf{ARM CMSIS-DSP} --- стандартизований пакет цифрової обробки
сигналів (ЦОС) для мікроконтролерів ARM Cortex-M. Надає жорстко
оптимізовані реалізації radix-4 для розмірів $N = 16, 64, 256, 1024,
4096$ із SIMD-оптимізацією NEON, але не переноситься на інші
архітектури~[9].

\textbf{NVIDIA CuFFT} та \textbf{OpenCL FFT} --- адаптації ШПФ для
графічних процесорів. Стандарт OpenCL забезпечує функціональну
портативність, проте не гарантує продуктивності: ядро, оптимізоване
під одну архітектуру GPU, демонструє низьку ефективність на іншій
через відмінності в ієрархії пам'яті та механізмах планування
потоків~[10].

\textbf{SPIRAL} (Carnegie Mellon University) --- генератор коду,
що створює платформо-залежну реалізацію безпосередньо з математичних
формул факторизації, описаних мовою SPL (Signal Processing Language).
Пошук оптимального графа зводиться до задачі пошуку у просторі
еквівалентних матричних розкладів~[11].

\textbf{Intel IPP} (Integrated Performance Primitives) --- жорстко
оптимізовані примітиви ШПФ під векторні інструкції SSE/AVX/AVX-512
процесорів Intel~[12].

\subsubsection*{Методи візуального моделювання}

\textbf{Загальні нотації} (UML~Activity, IDEF0/IDEF3, DFD
Йордана/ДеМарко) орієнтовані на control-flow парадигму, не мають
формальної семантики токенів і породжують надлишкові
\texttt{fork}/\texttt{join} конструкції для потокових обчислювачів.

\textbf{DFG/SDF} (Data Flow Graph і Synchronous Dataflow,
Лі та Месершмітт) --- спеціалізовані моделі для задач ЦОС.
У SDF-моделі кількість токенів, що споживаються та виробляються
кожним вузлом за одну активацію, є фіксованою і відомою на етапі
статичного аналізу, що дозволяє виконувати статичне планування без
ризику взаємних блокувань. Алгоритм ШПФ із фіксованою довжиною
$N = 16$ є природним прикладом SDF-графа~[5][6][7].

\textbf{Сигнальні графи і butterfly-діаграми} формалізують зв'язок
між матричною факторизацією та графом алгоритму~[4], але є
статичними і не моделюють динаміку виконання. У вітчизняній науковій
традиції аналогічний підхід застосовувався у роботах ОКБ~ЕІВТ
(1970-ті~рр.) та у Фізико-механічному інституті ім.~Г.~В.~Карпенка
НАН~України, де було розроблено векторний спецпроцесор ШПФ на основі
dataflow-парадигми~[2].

\subsubsection*{Порівняльний аналіз існуючих рішень}

Зведене порівняння розглянутих рішень за ключовими характеристиками
наведено у табл.~\ref{tab:solutions}.

\begin{table}[H]
  \caption{Порівняння існуючих рішень для обчислення ШПФ}
  \label{tab:solutions}
  \centering
  \small
  \begin{tabular}{p{2.4cm}p{2.0cm}p{1.8cm}p{1.6cm}p{1.4cm}p{2.6cm}}
    \toprule
    Рішення & Тип & Платформа & Динам. візуалізація & Відкр. доступ & Особливість \\
    \midrule
    FFTW 3.x      & SW бібл.        & CPU x86/ARM & Немає & Так   & Адаптивний planner \\
    CMSIS-DSP     & SW бібл.        & ARM Cortex-M & Немає & Так   & Фіксований radix-4 \\
    CuFFT         & SW бібл.        & GPU NVIDIA  & Немає & Ні    & Вендор-лок \\
    SPIRAL        & Генератор коду  & CPU/FPGA    & Немає & Так   & Пошук факторизацій \\
    Intel IPP     & SW бібл.        & Intel x86   & Немає & Ні    & AVX/SIMD \\
    Simulink (SDF)& Середовище     & Desktop     & Часткова & Ні & Комерційна ліцензія \\
    draw.io/Visio & Редактор діаграм & Desktop/Web & Ні    & Част. & Лише статика \\
    \textbf{Симулятор (БКР)} & \textbf{Веб-застос.} & \textbf{Браузер} & \textbf{Так} & \textbf{Так} & \textbf{DFG + верифікація} \\
    \bottomrule
  \end{tabular}
\end{table}

\subsubsection*{Висновок огляду}

Виявлено системну прогалину: жодне з розглянутих рішень не поєднує
одночасно достовірного числового обчислення ШПФ за формальною
SDF-моделлю, інтерактивної динамічної візуалізації руху токенів і
веб-доступності без комерційної ліцензії. Виробничі бібліотеки
(FFTW, CMSIS-DSP, CuFFT, SPIRAL, Intel~IPP) оптимізовані виключно
для продуктивності та не надають інтроспективних засобів. Теоретичні
моделі DFG/SDF формально визначені, проте рідко супроводжуються
веб-доступним симулятором із числовою верифікацією. Саме ця
незаповнена ніша визначає предмет розроблюваної БКР.

% ─────────────────────────────────────────────────
\subsection{Системний аналіз}

Для структурованого опису розриву між математичною моделлю,
формальним DFG/SDF-описом і виробничими реалізаціями ШПФ застосовано
системний аналіз: визначено зацікавлені сторони, ієрархію цілей,
обмеження та критерії успіху.

\subsubsection*{Зацікавлені сторони}

Потенційних користувачів і бенефіціарів симулятора об'єднано у
чотири групи (табл.~\ref{tab:stakeholders}).

\begin{table}[H]
  \caption{Зацікавлені сторони та їхні потреби}
  \label{tab:stakeholders}
  \centering
  \small
  \begin{tabular}{p{3.4cm}p{4.6cm}p{5.4cm}}
    \toprule
    Сторона & Потреба & Як симулятор її задовольняє \\
    \midrule
    Студенти спец.~122 «Комп'ютерні науки» &
      Зрозуміти роботу алгоритму ШПФ, побачити паралелізм і токенну динаміку &
      Анімація руху токенів, крок-за-кроком режим, теплова карта \\[2pt]
    Викладачі курсів ЦОС та паралельних обчислень &
      Демонстраційний інструмент для лекцій і лабораторних &
      Веб-застосунок без встановлення, відтворюваність на будь-якій ОС \\[2pt]
    Дослідники dataflow-архітектур &
      Перевірка гіпотез щодо планування та критичного шляху графа &
      Throughput/latency/queue, event log, налаштовувані затримки \\[2pt]
    Інженери DSP/FPGA &
      Попередня оцінка потокової моделі перед апаратною реалізацією &
      Параметричне налаштування \texttt{latency}/\texttt{delay}, експорт подій \\
    \bottomrule
  \end{tabular}
\end{table}

\subsubsection*{Ієрархія цілей}

\textbf{Генеральна ціль:} створити програмний симулятор, що відтворює
потокове виконання 16-точкового алгоритму ШПФ radix-4 $4\times4$ та
забезпечує його інтерактивну графічну візуалізацію у вигляді
анімованого DFG зі спостереженням токенів, метрик і наскрізної
затримки.

Підцілі, що випливають із генеральної:
\begin{enumerate}
  \item формалізувати граф потоку даних для radix-4 $4\times4$ з
        вузлами \texttt{source}, \texttt{dft4}, \texttt{twiddle},
        \texttt{sink} та дугами із параметричними затримками;
  \item реалізувати ядро симуляції з правилом активації токенів,
        FIFO-чергами та керованим логічним часом;
  \item забезпечити інтерактивну візуалізацію руху токенів,
        теплову карту активності та спектральний вихід;
  \item виконати числову верифікацію результатів шляхом порівняння
        з еталонним ДПФ ($\varepsilon \le 10^{-9}$).
\end{enumerate}

\subsubsection*{Обмеження}

\begin{itemize}
  \item \textit{технічні} --- цільова платформа: сучасний веб-браузер
        (Chrome/Firefox/Edge) з підтримкою JavaScript ES2020;
        відсутність серверної частини;
  \item \textit{ресурсні} --- розробка силами одного виконавця у
        межах БКР;
  \item \textit{часові} --- термін згідно з календарним графіком
        кафедри;
  \item \textit{алгоритмічні} --- фіксований розмір $N = 16$ як
        базовий демонстраційний; розширення на $N = 4^t$ ---
        напрямок подальшого розвитку;
  \item \textit{правові} --- лише програмні бібліотеки з відкритими
        ліцензіями (MIT, Apache~2.0, BSD).
\end{itemize}

\subsubsection*{Критерії успіху}

Результат вважається досягнутим, якщо одночасно виконуються такі
вимірювані критерії:
\begin{itemize}
  \item максимальне відхилення вихідного спектра від еталонного ДПФ
        не перевищує $\varepsilon = 10^{-9}$;
  \item чотири вузли \texttt{dft4} першого ярусу виконуються
        паралельно (одночасні події \texttt{onFire} у журналі);
  \item користувач може керувати швидкістю симуляції в діапазоні
        не менш ніж $0.1\times$--$10\times$ без втрати коректності;
  \item час відгуку UI на користувацькі дії не перевищує 100~мс.
\end{itemize}

\subsubsection*{Етичні та правові аспекти}

Симулятор працює з синтетичними тестовими сигналами, що генеруються
безпосередньо у браузері користувача. Проєкт не опрацьовує
персональних, медичних, біометричних чи фінансових даних, не передає
жодних даних на зовнішні сервери та не використовує моделей
машинного навчання для ухвалення рішень. Унаслідок цього вимоги
GDPR і Закону України «Про захист персональних даних» (№~2297-VI)
не поширюються на продукт у прямій формі. Робота відповідає
стандартам ISO/IEC~25010, IEEE~754, ECMAScript~2020,
ДСТУ~3008:2015 та ДСТУ~ГОСТ~7.1:2006.

% ─────────────────────────────────────────────────
\subsection{Розробка структури системи}

Архітектуру симулятора подано на двох рівнях моделі C4 (Context і
Container), доповнених деталізованим описом графа потоку даних
обчислювача та структур даних.

\subsubsection*{Контекстний рівень (C4 Level 1)}

На контекстному рівні система представляється як «чорна скринька»,
що взаємодіє з одним актором --- \emph{користувачем} (студент або
викладач) --- та з двома зовнішніми сутностями: \emph{браузером}
(runtime-середовище) і \emph{локальним сховищем браузера} (для
збереження пресетів вхідних сигналів і параметрів). Застосунок
завантажується одноразово зі статичного веб-хостингу та надалі
працює автономно, без серверної частини.

\subsubsection*{Рівень контейнерів (C4 Level 2)}

На рівні контейнерів симулятор декомпозується на чотири логічні
підсистеми, що виконуються у межах одного процесу браузера:

\begin{enumerate}
  \item \textbf{Simulation Core} --- ядро симуляції
        (\texttt{DataflowRuntime}, \texttt{scheduler}, \texttt{math});
        реалізує доставку токенів із FIFO-черг, перевірку правила
        активації та запуск готових вузлів; зберігає стан графа та
        черг.
  \item \textbf{Graph Generator} --- модуль побудови DFG для заданого
        $N$ та параметрів затримок (функція
        \texttt{generateDFT4x4()}).
  \item \textbf{State Store} --- централізоване сховище стану
        (паттерн \textit{single source of truth}), до якого
        підписуються всі UI-компоненти.
  \item \textbf{UI Layer} --- набір React-компонентів:
        \texttt{GraphView} (полотно React Flow), \texttt{SpectrumView}
        (діаграма $|X(k)|$), \texttt{Metrics} (throughput, latency,
        queue size), \texttt{Controls} (пуск/пауза/швидкість).
\end{enumerate}

Взаємодії між контейнерами: UI Layer читає стан зі State Store та
викликає команди; State Store делегує обробку до Simulation Core;
Simulation Core публікує події (\texttt{onToken}, \texttt{onFire},
\texttt{onOutput}), на які State Store оновлює відповідні поля;
Graph Generator викликається один раз при ініціалізації та за кожної
зміни $N$ чи \texttt{latency}.

\subsubsection*{Граф потоку даних обчислювача}

Детальна структура DFG для $N = 16$ складається з \textbf{56~вузлів}
і \textbf{64~дуг}, розподілених за п'ятьма ярусами:

\begin{enumerate}
  \item 16~вузлів \texttt{source} --- генерують токени відліків $x(n)$;
  \item 4~вузли \texttt{dft4} (рядкові) --- виконують стадію~1
        (рядкові 4-точкові ДПФ);
  \item 16~вузлів \texttt{twiddle} --- виконують стадію~2 (множення
        на $W_{16}^{r k_1}$);
  \item 4~вузли \texttt{dft4} (стовпцеві) --- виконують стадію~3;
  \item 16~вузлів \texttt{sink} --- збирають спектральні відліки $X(k)$.
\end{enumerate}

Затримка \texttt{latency} вузла \texttt{dft4} становить 400 умовних
одиниць, \texttt{twiddle} --- 200; затримка доставки кожної дуги
\texttt{delay} = 600.

\textbf{Критичний шлях} і \textbf{коефіцієнт паралелізму}.
Найдовший шлях від \texttt{source} до \texttt{sink}:
$T_{\mathrm{kr}} = 400 + 600 + 200 + 600 + 400 = 2200$ ум.~од.
Послідовний час виконання: $T_{\mathrm{ps}} = 8 \cdot 400 + 16 \cdot
200 = 6400$ ум.~од. Коефіцієнт паралелізму
$K_p = T_{\mathrm{ps}} / T_{\mathrm{kr}} \approx 2{,}9$, тобто
паралельне виконання незалежних вузлів забезпечує майже трикратне
прискорення порівняно з послідовним обчисленням.

\subsubsection*{Структури даних}

Інформаційне забезпечення симулятора визначається набором типів,
що є програмними відповідниками формальних об'єктів DFG:

\begin{itemize}
  \item \texttt{Complex} --- комплексне число у форматі
        \{\,\texttt{re}, \texttt{im}\,\}, IEEE~754 double
        ($\varepsilon_{\mathrm{mach}} \approx 2.22 \cdot 10^{-16}$);
  \item \texttt{Token} --- токен DFG з полями
        \texttt{id}, \texttt{value}, \texttt{t} (час надходження),
        \texttt{originT} (час народження у \texttt{source});
  \item \texttt{Edge} --- спрямована дуга
        (\texttt{from}/\texttt{to} = \{node, port\}, \texttt{delay},
        \texttt{label}, \texttt{labelValue});
  \item \texttt{NodeSpec} --- специфікація вузла
        (\texttt{id}, \texttt{kind}, \texttt{inPorts},
        \texttt{outPorts}, \texttt{latency}, \texttt{params});
  \item \texttt{Graph = \{nodes: NodeSpec[], edges: Edge[]\}} ---
        повний граф обчислювача (56~вузлів, 64~дуги).
\end{itemize}

\subsubsection*{Механізми обміну даними}

Оскільки симулятор не має серверної частини, координація підсистем
реалізується через \emph{внутрішню шину подій} і \emph{реактивне
сховище стану}. Ядро публікує три типи подій (\texttt{onToken},
\texttt{onFire}, \texttt{onOutput}), на які підписуються UI-компоненти
без прямих посилань на ядро (паттерн Publisher/Subscriber). Сховище
\texttt{simStore} містить групи полів: стан графа, стан виконання
(\texttt{IDLE}/\texttt{RUNNING}/\texttt{PAUSED}/\texttt{DONE}),
метрики (EMA throughput/latency/queue), результат (спектр,
event~log). Пресети вхідних сигналів зберігаються у Local~Storage
браузера.

\subsubsection*{Інтерфейс користувача}

Інтерфейс спроєктовано за принципом \emph{інформаційної ієрархії}:
найбільша площа екрана відведена під центральний об'єкт дослідження
--- граф потоку даних, --- а допоміжні панелі (метрики, спектр,
керування) розміщуються по периметру. Головний екран розділений на
чотири функціональні зони: центральна (\texttt{GraphView}), верхня
(\texttt{Controls}), права (\texttt{SpectrumView} + \texttt{Metrics}),
нижня (\texttt{TerminalOutput}).

Для підсилення сприйняття потокової природи алгоритму використовуються
три візуальні механізми: анімація токенів (кружечки з числовими
значеннями, що рухаються по дугах); теплова карта активності вузлів
(експоненціальне згасання $a_i(\tau + \mathrm{d}\tau) = 0{,}92 \cdot
a_i(\tau) + \Delta a_i$); кільцевий індикатор стану \texttt{busy} на
вузлах, що ілюструє затримку \texttt{latency}.

% ─────────────────────────────────────────────────
\subsection{Вибір технологій реалізації системи}

Технології реалізації обрано на основі трирівневого аналізу
альтернатив: архітектурний підхід, парадигма реалізації UI та стану,
метод візуалізації DFG.

\subsubsection*{Архітектурний підхід}

Розглянуто три варіанти:

\textbf{Варіант~А: настільний застосунок} (Electron, Qt, JavaFX).
Забезпечує повний доступ до ресурсів хост-системи та багатопоточність,
проте вимагає встановлення на кожну робочу станцію та різних збірок
під Windows, Linux, macOS, що суперечить вимогам NF7 (без встановлення)
і ускладнює дотримання NF4 (портативність). \emph{Відхилено}.

\textbf{Варіант~Б: клієнт-серверний веб-застосунок.} Тонкий клієнт у
браузері виконує лише візуалізацію, а ядро розміщується на сервері.
Вносить мережеву затримку (несумісно з NF2 $\le 100$~мс), потребує
підтримки інфраструктури та ускладнює відтворюваність. \emph{Відхилено}.

\textbf{Варіант~В: повністю клієнтський SPA} (Single-Page Application).
Усі обчислення на стороні браузера; серверної частини немає.
Завантаження --- одноразовий статичний пакет (HTML+JS+CSS); надалі
робота не потребує мережі. Одночасно задовольняє NF2, NF4, NF5, NF7.
\emph{Обрано}.

\subsubsection*{Парадигма реалізації UI та стану}

Для координації трьох принципово різних підсистем (графова
візуалізація, числове ядро, панель керування) обрано поєднання
\textbf{реактивного рендерингу} (React) і \textbf{централізованого
спостерігаваного сховища стану} (\texttt{Zustand}). Реактивний
рендеринг автоматично перемальовує лише ті компоненти, що підписані
на змінене поле сховища (NF3, 60~fps). Централізоване сховище
гарантує, що всі панелі бачать узгоджений «знімок» стану.
Альтернативи: пряме маніпулювання DOM (jQuery-style) --- лінійне
зростання складності; Flux/Redux --- надмірний шаблонний код для
базових сценаріїв.

\subsubsection*{Метод візуалізації DFG}

Розглянуто три підходи:
\begin{itemize}
  \item \textbf{HTML-canvas} --- максимальна продуктивність, але
        ручна реалізація перетягування, виділення та підказок;
  \item \textbf{SVG + ручне кодування} --- простий у налагодженні,
        але обмежено розширюваний;
  \item \textbf{React~Flow} --- спеціалізована бібліотека з готовими
        механізмами перетягування, зума, анімованих ребер та
        формальною моделлю \texttt{Node}/\texttt{Edge}, що добре
        відображається на абстракцію DFG. \emph{Обрано}.
\end{itemize}

\subsubsection*{Шаблони проєктування}

У системі застосовано класичні шаблони:
\begin{itemize}
  \item \textit{Observer (Publisher/Subscriber)} --- події
        \texttt{onToken}, \texttt{onFire}, \texttt{onOutput};
  \item \textit{State / Finite State Machine} --- життєвий цикл
        симуляції \texttt{IDLE} $\to$ \texttt{RUNNING} $\to$
        \texttt{PAUSED} $\to$ \texttt{DONE};
  \item \textit{Strategy} --- різні типи вузлів за полем
        \texttt{NodeKind};
  \item \textit{Factory} --- генерація графа функцією
        \texttt{generateDFT4x4()};
  \item \textit{Facade} --- клас \texttt{DataflowRuntime} приховує
        реалізацію черг і планувальника за єдиним методом
        \texttt{tick(dt, maxFires)}.
\end{itemize}

\subsubsection*{Типізація та забезпечення якості}

Для дотримання NF1 ($\varepsilon \le 10^{-9}$) і NF5 (детермінованість)
обрано \textbf{TypeScript} (Apache~2.0) як мову реалізації зі
статичною типізацією. Це унеможливлює помилки передачі аргументів
неправильного типу на етапі компіляції. Комплексні числа подаються
у форматі IEEE~754 double, що забезпечує запас точності відносно
$\varepsilon$.

\subsubsection*{Стек технологій}

Підсумковий стек технологій реалізації наведено у
табл.~\ref{tab:tech_stack}.

\begin{table}[H]
  \caption{Стек технологій реалізації симулятора}
  \label{tab:tech_stack}
  \centering
  \small
  \begin{tabular}{p{3.0cm}p{3.0cm}p{2.0cm}p{5.5cm}}
    \toprule
    Шар & Технологія & Ліцензія & Призначення \\
    \midrule
    Мова          & TypeScript    & Apache 2.0 & Статична типізація, безпека типів \\
    UI-фреймворк  & React 18      & MIT        & Реактивний рендеринг компонентів \\
    Метафреймворк & Next.js 14    & MIT        & Збірка SPA, статичний експорт \\
    Сховище стану & Zustand       & MIT        & Централізоване спостерігаване сховище \\
    Візуалізація  & React Flow    & MIT        & Граф потоку даних із готовою інтерактивністю \\
    Стилі         & SCSS modules  & MIT        & Локальні стилі компонентів \\
    Тестування    & Vitest        & MIT        & Юніт-тести ядра і структур даних \\
    Верифікація   & fft.js        & MIT        & Еталонна реалізація для \texttt{compareWithFFT()} \\
    \bottomrule
  \end{tabular}
\end{table}

Зв'язок між нефункціональними вимогами та обраними технологічними
рішеннями зведено у табл.~\ref{tab:req_to_tech}.

\begin{table}[H]
  \caption{Зв'язок вимог і архітектурних/технологічних рішень}
  \label{tab:req_to_tech}
  \centering
  \small
  \begin{tabular}{p{2.2cm}p{5.6cm}p{6.0cm}}
    \toprule
    Вимога & Рішення & Обґрунтування \\
    \midrule
    NF1 (точність)        & TypeScript + IEEE 754 double         & Виключення помилок типів, $\varepsilon_{\mathrm{mach}}{\approx}2.22{\cdot}10^{-16}$ \\
    NF2 (відгук $\le 100$~мс) & Клієнтський SPA без серверних викликів & Відсутність мережевої затримки \\
    NF3 (60~fps)          & React + Zustand із селекторами       & Перемальовування лише змінених компонентів \\
    NF4 (портативність)   & ES2020, Web APIs                     & Робота у Chrome/Firefox/Edge без polyfill \\
    NF5 (детермінованість) & Логічний час + статичний граф       & Ідентичні входи $\Rightarrow$ ідентичний журнал \\
    NF6 (відкритий код)   & Тільки OSS-бібліотеки                & React, React Flow, Zustand, TypeScript \\
    NF7 (без встановлення) & SPA у статичному хостингу           & Одноразове завантаження бандлу \\
    NF8 (зручність)       & React Flow + панелі Controls/Metrics & Двокліковий доступ до базових операцій \\
    \bottomrule
  \end{tabular}
\end{table}

% =============================================================
\subsection*{Висновки}

У межах практики виконано підготовчий етап бакалаврської
кваліфікаційної роботи за темою «Візуальне моделювання роботи
обчислювача ШПФ під керуванням потоку даних». Зокрема:

\begin{enumerate}
  \item сформульовано постановку задачі та визначено перелік із
        9~функціональних (F1--F9) і 8~нефункціональних
        (NF1--NF8) вимог; розроблено перелік із 9~прецедентів
        використання, що охоплюють основні сценарії роботи студента
        і викладача;
  \item проведено огляд літератури та аналогів: проаналізовано
        виробничі бібліотеки (FFTW, ARM CMSIS-DSP, NVIDIA CuFFT,
        SPIRAL, Intel IPP) і методи візуального моделювання
        (UML, IDEF, DFG/SDF, сигнальні графи); виявлено системну
        прогалину --- відсутність веб-доступного симулятора, що
        поєднує формальну SDF-модель, виконуваний алгоритм із
        числовою верифікацією та інтерактивну візуалізацію;
  \item виконано системний аналіз: визначено зацікавлені сторони
        (студенти, викладачі, дослідники, інженери DSP/FPGA),
        ієрархію цілей, обмеження (технічні, ресурсні, часові,
        алгоритмічні, правові) та вимірювані критерії успіху
        ($\varepsilon \le 10^{-9}$, відгук $\le 100$~мс, 60~fps,
        паралельне виконання вузлів);
  \item розроблено структуру системи: архітектура C4 Context/Container
        з декомпозицією на чотири підсистеми (Simulation Core, Graph
        Generator, State Store, UI Layer); деталізовано граф
        обчислювача (56~вузлів, 64~дуги, $T_{\mathrm{kr}} = 2200$,
        $K_p \approx 2{,}9$); визначено структури даних
        (\texttt{Complex}, \texttt{Token}, \texttt{Edge},
        \texttt{NodeSpec}, \texttt{Graph}); описано механізми обміну
        даними та інтерфейс користувача;
  \item обґрунтовано вибір технологій реалізації: клієнтський SPA на
        TypeScript + React + Next.js + Zustand + React Flow з
        еталонною бібліотекою fft.js для верифікації; усі обрані
        компоненти мають відкриті ліцензії (MIT, Apache~2.0), що
        задовольняє NF6.
\end{enumerate}

Отримані результати утворюють повне технічне завдання для
безпосередньої реалізації симулятора, що складатиме практичну
частину БКР.

% =============================================================
\section*{Список використаних джерел}
\addcontentsline{toc}{section}{Список використаних джерел}

\begin{enumerate}
  \item Nussbaumer~H.~J. Fast Fourier Transform and Convolution
        Algorithms. --- 2nd corr. and updated ed. --- Vol.~2. ---
        Springer Series in Information Sciences. --- Berlin~;
        Heidelberg : Springer-Verlag, 1982. --- 276~p.

  \item Процько~І.~О., Теслюк~В.~М., Дунець~Р.~Б. Швидкому
        перетворенню Фур'є --- 60~років // Ukrainian Journal of
        Information Technology. --- 2025. --- Vol.~7, No.~1. --- P.~1--1.

  \item Cooley~J.~W., Tukey~J.~W. An Algorithm for the Machine
        Calculation of Complex Fourier Series // Mathematics of
        Computation. --- 1965. --- Vol.~19, No.~90. --- P.~297--301. ---
        \url{https://doi.org/10.1090/S0025-5718-1965-0178586-1}.

  \item Burrus~C.~S., Parks~T.~W. DFT/FFT and Convolution Algorithms:
        Theory and Implementation. --- New York : John Wiley \& Sons,
        1985. --- 240~p.

  \item Dennis~J.~B. First Version of a Data Flow Procedure Language //
        Programming Symposium. Lecture Notes in Computer Science. ---
        1974. --- Vol.~19. --- P.~362--376. ---
        \url{https://doi.org/10.1007/3-540-06859-7_145}.

  \item Dennis~J.~B., Misunas~D.~P. A Preliminary Architecture for a
        Basic Data-Flow Processor // Proceedings of the 2nd Annual
        Symposium on Computer Architecture. --- 1975. --- P.~126--132.

  \item Veen~A.~H. Dataflow Machine Architecture // ACM Computing
        Surveys. --- 1986. --- Vol.~18, No.~4. --- P.~365--396. ---
        \url{https://doi.org/10.1145/27633.28055}.

  \item Frigo~M., Johnson~S.~G. The Design and Implementation of
        FFTW3 // Proceedings of the IEEE. --- 2005. --- Vol.~93,
        No.~2. --- P.~216--231. ---
        \url{https://doi.org/10.1109/JPROC.2004.840301}.

  \item ARM Limited. CMSIS-DSP Software Library : Technical Reference
        Manual. --- Cambridge : ARM Limited, 2023.

  \item Du~P., Weber~R., Luszczek~P., Tomov~S., Peterson~G.,
        Dongarra~J. From CUDA to OpenCL: Towards a Performance-Portable
        Solution for Multi-Platform GPU Programming // Parallel
        Computing. --- 2012. --- Vol.~38, No.~8. --- P.~391--407. ---
        \url{https://doi.org/10.1016/j.parco.2011.10.002}.

  \item P\"uschel~M., Moura~J.~M.~F., Singer~B. et al. SPIRAL: A
        Generator for Platform-Adapted Libraries of Signal Processing
        Algorithms // International Journal of High Performance
        Computing Applications. --- 2004. --- Vol.~18, No.~1. ---
        P.~21--45. ---
        \url{https://doi.org/10.1177/1094342004041291}.

  \item Intel Corporation. Intel Integrated Performance Primitives
        (Intel IPP) : Developer Guide and Reference. --- Santa Clara :
        Intel Corporation, 2021.
\end{enumerate}

\end{document}
